\documentclass{article}
\usepackage{amsmath}
\usepackage{graphicx}
\graphicspath{ {images/} }
\usepackage[spanish]{babel}
\usepackage{bbm}
\usepackage{amsthm}
\usepackage{authblk}
\usepackage{hyperref}
\usepackage[utf8]{inputenc}
\usepackage{mathrsfs}
\usepackage{amsfonts}
\usepackage{amssymb}
\usepackage{enumerate}
\usepackage[left=3cm,right=2cm,top=2cm,bottom=2cm]{geometry}
\usepackage{epsfig}
 \renewcommand{\baselinestretch}{0.93}
 \thispagestyle{empty}
 \pagestyle{empty}
\setlength{\textheight}{53pc} \setlength{\textwidth} {36pc}
\setlength{\topmargin}{-4mm}
\usepackage{multicol}
\newcommand*{\email}[1]{%
    \normalsize\href{mailto:#1}{#1}\par
    }
\setlength{\columnsep}{1cm}
\newcommand{\N}{\mathbb{N}}
\newcommand{\calM}{\cal{M}}
\newcommand{\calP}{\cal{P}}
\newcommand{\F}{\mathbb{F}}
\newcommand{\R}{\mathbb{R}}
\newcommand{\Z}{\mathbb{Z}}
\newcommand{\Q}{\mathbb{Q}}
\newcommand{\C}{\mathbb{C}}
\newcommand{\bbH}{\mathbb{H}}


\theoremstyle{definition}
\newtheorem{ejer}{Ejercicio}
\author{Marcos Morales}
\affil{Universidad Finis Terrae}
\title{Primera Semana\\}



	


\begin{document}


\begin{picture}(400, 75)
   \put(-40,15){\includegraphics[width=5cm]{UFT.png}}
\end{picture}


\begin{center}
\huge{Octava Semana}
\end{center}

\begin{center}
\Large{ Marcos Morales, Camila Muñoz}
\end{center}

\begin{center}
	\large{Cálculo Vectorial}
\end{center}
\begin{center}
\setlength{\parindent}{0cm}\large{Clases centradas para capacitar a los estudiantes de ingeniería en Cálculo Vectorial. \\}
\end{center}


\vspace{1cm}

\begin{center}
	\large{Contenidos:} 
\end{center}
\vspace{0.1cm}
\begin{center}
\begin{minipage}{0.5\textwidth}
\begin{enumerate}
    \item Cambio de Variable
    \item Coordenadas Polares
    \item Aplicaciones de las integrales dobles
\end{enumerate}
\end{minipage}
\end{center}




\vspace{6cm}




\pagestyle{empty}


\setcounter{page}{1}
\pagestyle{plain}
\setlength{\parindent}{0cm}





\newpage

\section{Cambio de variable}

Cuando ten\'iamos que realizar un cambio de variable en integrales simples digamos la integral

\begin{align*}
    \int_a^b f(x)dx
\end{align*}

Con el cambio de variable $x=g(u)$, entonces la nueva integral quedaba


\begin{align*}
    \int_{g(a)}^{g(b)} f(g(u))g'(u)du
\end{align*}

Esto es lo que se conoce como cambio de variable y solamente se usa la regla de la cadena. Podemos ver que cambian los l\'imites de integraci\'on y adem\'as aparece $g'(u)$ multiplicando. \\


En caso de que tengamos dos variables, necesitamos reemplazarlas por otras dos variables. Por ejemplo $x=g(u,v)$ e $y=h(u,v)$, donde $u,v$ ser\'an las nuevas variables. Naturalmente, si la regi\'on encerrada por $x$ e $y$ es $R_{xy}$ en la integral doble original, entonces la nueva regi\'on ser\'a la delimitada por $u$ y $v$ , que llamaremos $R_{u,v}$, de la misma forma como cambian los l\'imites en integrales simples. 

\begin{figure}[h]
    \centering
    \includegraphics[width=0.8\textwidth]{Imagenes/cambio.png}
    \caption{Cambio de regi\'on al hacer un cambio de variable}
    \label{fig:etiqueta}
\end{figure}


Pero la pregunta es, qu\'e reemplaza a $g'(u)$?, para responder esta pregunta consideramos la transformaci\'on lineal $T$, como $T(u,v)=(x(u,v),y(u,v))$, entonces 

\begin{align*}
    J(T) = \frac{\partial(x,y)}{\partial(u,v)} = \begin{vmatrix}
        \displaystyle\frac{\partial x}{\partial u} &  \displaystyle\frac{\partial x}{\partial v} \\
        \displaystyle\frac{\partial y}{\partial u} &  \displaystyle\frac{\partial y}{\partial v} 
    \end{vmatrix}
\end{align*}

este determinante es llamado el jacobiano de cambio de base, el valor absoluto del jacobiano es el que reemplaza a la derivada dentro de la integral.\\

\textbf{Teorema (Cambio de variable):} Sea $T : R_{uv} \subset \mathbb{R}^2 \to R_{xy}\subset \mathbb{R}^2$ una transformaci\'on inyectiva y de clase $C^1$ en el plano $uv$ en el plano $xy$ tal que $T (u, v ) = (x (u, v ) , y (u, v ))$. Si $f : R_{xy} \to \mathbb{R}$ es Riemann integrable, entonces:

\begin{align*}
    \iint_{R_{xy}} f(x,y)dA =  \iint_{R_{u,v}} f(u,v) |J(T)|dA
\end{align*}

\newpage

\textbf{Ejercicio 1:} Calcule el valor de $\iint_{R} x+y+1dA$ donde $R$ es la regi\'on acotada por las gr\'aficas de $y-x=1$, $y-x=-1$, $y+x=1$ e $y+x=2$ \\

\textbf{Respuesta:} Este ejercicio es muy dif\'icil de hacer sin usar cambio de variable, por suerte podemos usar el cambio de variable para hacerlo de forma sencilla. \\

El cambio que hay que realizar es $u=y-x$ y $v=y+x$, ahora tenemos que despejar $x$ e $y$ en funci\'on de $u$ y $v$, sumando ambos obtenmos que $y$ y restando obtenemos $x$

\begin{align*}
    y &= \frac{u+v}{2} \\
    x &= \frac{v-u}{2}
\end{align*}

Luego


\begin{align*}
    J(T) &=\begin{vmatrix}
        -\frac{1}{2} & \frac{1}{2} \\
        \frac{1}{2} & \frac{1}{2} \\
    \end{vmatrix} \\
    &= -\frac{1}{2}
\end{align*}

Luego $|J(T)| =1/2$, los l\'imites para laa variable $u$ son $-1$ y $1$ mientras que para la variable $v$ son $1$ y $2$, esto es claro debido a la definici\'on de la regi\'on de integraci\'on, por lo tanto




\begin{align*}
    \iint_{R_{xy}} f(x,y)dA &=  \iint_{R_{u,v}} f(T(u,v)) |J(T)|dA \\
    &=  \frac{1}{2}\int_{-1}^1\int_1^2  \frac{v-u}{2} +\frac{u+v}{2} +1 dvdu \\
    &=  \frac{1}{2}\int_{-1}^1\int_1^2  v +1 dvdu \\
    &= \frac{1}{2} \int_{-1}^1\left( \frac{v^2}{2}+v\right)\bigg|_1^2 du \\
    &= \frac{1}{2} \int_{-1}^1 \frac{5}{2}du \\
    &= \frac{1}{2} \cdot \frac{5}{2} \cdot 2 \\
    &= \frac{5}{2}
\end{align*}


\textbf{Ejercicio 2:} Calcule el valor de $\displaystyle \iint_{R} \sqrt{\frac{x}{y}}+ \sqrt{xy} dA$ donde $R$ es la regi\'on acotada por las gr\'aficas de $xy=1$, $yx=9$, $y=x$ e $y=4x$ \\

\textbf{Respuesta:} Es claro que uno de los cambios de variable que queremos hacer es $u=xy$, el otro cambio no est\'a tan claro, pero si dividimos por que las segundas dos condiciones, obtenemos que $y/x= 1$ e $y/x=4$, por lo tanto el segundo cambio de variable conveniente es $v= y/x$. \\

Despejando $y$ en ambas igualdades tenemos $y=u/x=vx$, luego $x^2=\displaystyle \frac{u}{v}$ y por lo tanto $\displaystyle x =\sqrt{\frac{u}{v}}$, luego $y = \displaystyle vx = \sqrt{uv}$. Entonces 

\begin{align*}
    \frac{\partial x}{\partial u} = \frac{1}{2\sqrt{uv}} \hspace{0.3cm}   \frac{\partial x}{\partial v} = -\frac{\sqrt{u}}{\sqrt{v^3}} 
\end{align*}

\begin{align*}
    \frac{\partial y}{\partial u} = \frac{\sqrt{}}{2} \hspace{0.3cm}   \frac{\partial x}{\partial v} = -\frac{\sqrt{u}}{\sqrt{v^3}} 
\end{align*}



De esta forma el jacobiano queda

\begin{align*}
|J| = \left| \frac{\partial(x, y)}{\partial(u, v)} \right| 
= \left|
\begin{vmatrix}
\frac{1}{2} u^{-1/2} v^{-1/2} & -\frac{1}{2} u^{1/2} v^{-3/2} \\
\frac{1}{2} u^{-1/2} v^{1/2} & \frac{1}{2} u^{1/2} v^{-1/2}
\end{vmatrix}
\right|
= \frac{1}{2v}.
\end{align*}

\vspace{0.3cm}

Los nuevos l\'imites var\'ian entre

\begin{align*}
1 \leq u \leq 9, \\
1 \leq v \leq 4. 
\end{align*}

La nueva funci\'on a integrar ser\'a

\begin{align*}
\sqrt{\frac{x}{y}} + \sqrt{xy} = \sqrt{\frac{1}{v}} + \sqrt{u}.
\end{align*}


\vspace{0.3cm}

De esta fomra la integral queda


\begin{align*}
    \iint_R \left( \sqrt{\frac{x}{y}} + \sqrt{xy} \right) dA 
&= \iint_{R_{uv}} \left( \sqrt{\frac{1}{v}} + \sqrt{u} \right) \cdot \frac{1}{2v} \, du\,dv \\
&=  \int_{1}^{9} \int_{1}^{4} \left( \frac{v^{-3/2}}{2} + \frac{u^{1/2}}{2v} \right) dv\, du \\
&=\int_{1}^{9} \left( -v^{-1/2} + \frac{u^{1/2}}{2}\ln(v) \right)\bigg|_1^4 du\\
&=\int_{1}^{9} -\frac{1}{2} + \frac{u^{1/2}}{2}\ln(4)+1 du \\
&=\int_{1}^{9} \frac{1}{2}+ \frac{u^{1/2}}{2}\ln(4) du \\
&= \left(\frac{u}{2}+ \frac{u^{3/2}}{3}\ln(4) \right)\bigg|_1^9 \\
&= 4+\frac{26\ln(4)}{3}
\end{align*}



\newpage

\section{Coordenadas Polares}

Reecoredemos que todo punto $P(x,y)$ en el plano cartesiano puede ser escrito en coordenadas polares

\begin{figure}[h]
    \centering
    \includegraphics[width=0.5\textwidth]{Imagenes/polares.png}
    \caption{Punto en coordenadas polares}
    \label{fig:etiqueta}
\end{figure}

El cambio de variable es

\begin{align*}
    x &= r\cos(\theta) \\
    y &= r\sin(\theta) \\
    r &= \sqrt{x^2+y^2} \\
    \theta &= \arctan\left( \frac{y}{x}\right)
\end{align*}

Esto nos permite transformar regiones $R_{xy}$ en el plano $xy$ a una regi\'on $R_{r\theta}$ en el plano $r\theta$

\begin{figure}[h]
    \centering
    \includegraphics[width=0.3\textwidth]{Imagenes/cambiopolar.png}
    \caption{Regi\'on en coordenadas polares}
    \label{fig:etiqueta}
\end{figure}
\begin{align*}
    R_{r\theta} = \lbrace (r,\theta) : \alpha \leq \theta \leq \beta, h_1(\theta) \leq r \leq h_2(\theta) \rbrace
\end{align*}

Por lo tanto consideramos la transformaci\'on lineal $T: \mathbb{R}^2 \to \mathbb{R}^2$ dada por $T(r,\theta)=(r\cos(\theta),r\sin(\theta))$. Esta transformaci\'on es inyectiva para $r \geq 0$ y $0 < \theta \leq 2\pi$ o bien  $0 \leq \theta < 2\pi$. Su jacobiano es

\begin{align*}
    J(T) &= \begin{vmatrix}
        \cos(\theta) & -r\sin(\theta) \\
        \sin(\theta) & r\cos(\theta) 
    \end{vmatrix} \\
    &=r
\end{align*}

por lo tanto usando el teorema de cambio de variable obtenemos el siguiente resultado


\begin{align*}
    \iint_{R_{xy}} f(x,y)dA = \iint_{R_{r,\theta}} f(r\cos(\theta),r\sin(\theta))rdA
\end{align*}

\newpage



\textbf{Ejercicio 1:} Calcular 

\begin{align*}
    \int_0^1 \int_0^{\sqrt{1-x^2}} e^{\sqrt{x^2+y^2}}dydx
\end{align*}

\textbf{Respuesta:} Aqu\'i nos conviene pasar a coordenadas polares, primero hay que entender la regi\'on de integraci\'on, tenemos que $x$ var\'ia entre 0 y 1 mientras que $y$ var\'ia entre la recta $y=0$ y la curva $y= \sqrt{1-x^2}$, esto corresponde a un cuarto de circunferencia de radio 1

\begin{figure}[h]
    \centering
    \includegraphics[width=0.3\textwidth]{Imagenes/cuartocircunferencia.png}
    \caption{Regi\'on en coordenadas polares}
    \label{fig:etiqueta}
\end{figure}

Podemos ver que $0 \leq \theta < \pi/2$ y $0 \leq r \leq 1$, adem\'as sabemos que $r=\sqrt{x^2+y^2}$, luego


\begin{align*}
    \int_0^1 \int_0^{\sqrt{1-x^2}} e^{\sqrt{x^2+y^2}}dydx &= 
    \int_0^{\pi/2}\int_0^1 re^r dr d\theta \\
    &= \int_0^{\pi/2} (re^r-e^r)\bigg|_0^1 d\theta \\
    &= \int_0^{\pi/2} 1 d\theta \\
    &= \frac{\pi}{2}
\end{align*}



\textbf{Ejercicio 2:} Calcular 

\begin{align*}
    \int_{0}^3 \int_0^{\sqrt{9-x^2}} \arctan\left( \frac{y}{x}\right)dydx
\end{align*}

\textbf{Respuesta:} Tenemos que $x$ var\'ia entre 0 y 3 mientras que $y$ var\'ia entre la recta $y=0$ y la curva $y= \sqrt{9-x^2}$, esto corresponde a un cuarto de circunferencia de radio 3. \\



Podemos ver que $0 \leq \theta < \pi/2$ y $0 \leq r \leq 3$, adem\'as sabemos que $r=\arctan(y/x)$, luego


\begin{align*}
    \ \int_{-3}^3 \int_0^{\sqrt{9-x^2}} \arctan\left( \frac{y}{x}\right)dydx &= 
    \int_0^{\pi/2}\int_0^3 \theta r dr d\theta \\
    &= \int_0^{\pi/2} \theta \left( \frac{r^2}{2}\right)\bigg|_0^3 d\theta \\
    &= \frac{9}{2}\int_0^{\pi/2} \theta d\theta \\
    &= \frac{9}{2}\left( \frac{\theta^2}{2}\right)\bigg|_0^{\pi/2} d\theta \\
    &= \frac{9\pi^2}{16}
\end{align*}

\newpage

\section{Aplicaciones de las integrales dobles}

\subsection{Centro de masas}

Las integrales dobles tienen muchas aplicaciones, una de ellas es encontrar el centro de masas de. Supongamos que una l\'amina ocupa una rgi\'on $R$ del plano $xy$ y su densidad en un punto $(x,y)$ est\'a dada por $\rho(x,y)$ , donde $\rho $ es una funci\'on continua sobre $D$. La masa total $m$ de la l\'amina es

\begin{align*}
    m &= \iint_D \rho(x,y) dA
\end{align*}

Las coordenadas $(\overline{x},\overline{y})$ del centro de masas vienen dadas por

\begin{align*}
    \overline{x} &= \frac{1}{m} \iint_D x\rho(x,y) dA \\
    \overline{y} &= \frac{1}{m} \iint_D y\rho(x,y) dA \\
\end{align*}

\textbf{Ejercicio:} Determine la masa y el centro de masa  de la l\'amina que ocupa la regi\'on $D$ del primer cuadrante  acotada por $y=x^2$ y la recta $y=1$ y que tiene funci\'on de densidad $\rho(x,y) = xy$ \\

\textbf{Respuesta: } La masa ser\'a

\begin{align*}
    m &= \int_{0}^1\int_{x^2}^1 xy dydx \\
      &= \int_{0}^1 \frac{xy^2}{2}\bigg|_1^{x^2}dx \\
      &= \int_{0}^1 \frac{x}{2}-\frac{x^5}{2} dx \\
      &=\left(\frac{x^2}{4}-\frac{x^6}{12}\right)\bigg|_{0}^1  \\
      &= \frac{1}{6}
\end{align*}

Luego

\begin{align*}
    \overline{x} &= \frac{1}{m} \int_{0}^1\int_{x^2}^1 x^2y dydx \\
        &= 6 \int_{0}^1 \frac{x^2y^2}{2}\bigg|_{x^2}^1 dx \\
        &= 6 \int_{0}^1 \frac{x^2}{2}-\frac{x^6}{2} dx \\
        &= 6 \left( \frac{x^3}{6}-\frac{x^7}{14} \right)\bigg|_0^1\\
        &=\frac{4}{7}
\end{align*}

Por otra parte

\begin{align*}
    \overline{y} &= \frac{1}{m} \int_{0}^1\int_{x^2}^1 xy^2 dydx \\
        &= 6 \int_{0}^1 \frac{xy^3}{3}\bigg|_{x^2}^1 dx \\
        &= 6 \int_{0}^1 \frac{x}{3}-\frac{x^7}{3} dx \\
        &= 6 \left( \frac{x^2}{6}-\frac{x^8}{24} \right)\bigg|_0^1\\
        &=\frac{3}{4}
\end{align*}

Por lo tanto la masa total es $1/6$ y el centro de masas es $C=\displaystyle \left( \frac{4}{7},\frac{3}{4}\right)$




\newpage


\subsection{Momento de Inercia}

Supongamos que una l\'amina ocupa una rgi\'on $R$ del plano $xy$ y su densidad en un punto $(x,y)$ est\'a dada por $\rho(x,y)$, entonces las integrales


\begin{align*}
    I_x &= \iint_R y^2\rho(x,y) dA \\
    I_y &= \iint_R x^2\rho(x,y) dA \\
\end{align*}

Corresponden a los momentos de inercia de la l\'amina al rededor del eje $x$ y el eje $y$ respectivamente. El momento de inercia al rededor del origen, tambi\'en llamado momento polar es

\begin{align*}
I_0 =\iint_D (x^2+y^2)\rho(x,y)dA
\end{align*}

\textbf{Ejercicio:} Determine el momento de inercia del eje $X$ de la l\'amina correspondiente a la regi\'on parab\'olica en el primer cuadrante dada por  $0 \leq y \leq 4-x^2$ si $\rho(x,y) = xy$ \\

\textbf{Respuesta: } El momento de inercia pedido viene dado por

\begin{align*}
    I_x &= \int_{0}^2 \int_0^{4-x^2} xy^3 dy dx \\
    &= \int_{0}^2 \left(\frac{xy^4}{4}\right)\bigg|_0^{4-x^2} dx \\
    &= \int_{0}^2 x\frac{(4-x^2)^4}{4} dx \\
\end{align*}

Para resolver esta integral usamos el cambio de variable $u=4-x^2$ entonces $du=-2xdx$, cuando $x=0$ se tiene $u=4$ y caundo $x=2$ se tien $u=0$, luego


\begin{align*}
    I_x &= \int_{0}^2 x\frac{(4-x^2)^4}{4} dx \\
    &= \int_{4}^0 -\frac{u^4}{8} du \\
    &= -\frac{u^5}{40}\bigg|_4^0 \\
    &=\frac{4^5}{40} \\
    &= \frac{128}{5}
\end{align*}


\newpage


\subsection{Área de una superficie}

Consideremos la superficie $S$ dada por $z=f(x,y)$ definida sobre una regi\'on  $D$ cerrada y acotada del plano $XY$. Si $f$ y sus primeras derivadas parciales son continuas en la regi\'on $D$, entonces el \'area de la superficie $S$ es

\begin{align*}
    A(S) = \iint_{D} \sqrt{1-(f_x(x,y))^2-(f_y(x,y))^2} dA
\end{align*}


\textbf{Ejemplo: } Calcular el \'area de la proci\'on del plano $z=2-x-y$ que esta situada encima del ci\'rculo $x^2+y^2 \leq 1$ en el primer cuadrante \\

\textbf{Respuesta: } Dado que \( z = f(x,y)= 2 - x - y \), tenemos

\[
f_x = -1, \quad f_y = -1
\]

Por lo tanto, el integrando se convierte en

\[
\sqrt{1 + (-1)^2 + (-1)^2} = \sqrt{3}
\]

La región \( D \) es el cuarto de círculo de radio 1 en el primer cuadrante, es decir

\[
x^2 + y^2 \leq 1, \quad x \geq 0, \quad y \geq 0
\]

Usamos coordenadas polares

\[
x = r \cos \theta, \quad y = r \sin \theta
\]

Los límites de integración son

\[
0 \leq r \leq 1, \quad 0 \leq \theta \leq \frac{\pi}{2}
\]

Entonces el área es

\begin{align*}
    A &=  \iint_D \sqrt{3} \, dx\,dy \\
    &=  \sqrt{3} \int_0^{\frac{\pi}{2}} \int_0^1 r\,dr\,d\theta \\
    &=  \sqrt{3} \int_0^{\frac{\pi}{2}}  \frac{r^2}{2} \bigg|_0^1 \, d\theta \\
    &= \sqrt{3} \int_0^{\frac{\pi}{2}} \frac{1}{2} \, d\theta \\
    &= \frac{\pi \sqrt{3}}{4}
\end{align*}




\end{document} 